\section{기록 영상의 고정 모델 오프라인 사례 연구}\label{sec:case}
본 사례는 리프터 운용 중 저장된 네 RGB 영상 8,910프레임에 Base와 Base+N3를 각각 한 번 적용하여 코너·자세의 연속 출력을 비교하는 오프라인 평가이다. 촬영 당시 특정 팔레트에 맞춰 학습한 운용 모델과 논문 평가 모델을 구분하고, 두 평가 경로의 전처리·카메라·치수·좌표계를 고정하였다. 시간축 추적을 포함한 팔레트 연구\cite{mohamed2020}와 달리 새로운 추적기나 제어기를 제안하지 않는다. 두 방법은 각각 8,772프레임에서 새 자세를 산출했고, 이전 출력 유지 없이 138프레임이 no-pose였다.

CSV의 운용 모델 추정값이나 제어 명령은 독립 정답이 아니다. 영상 시각과의 대응, 물체 원점·축·단위를 확인한 뒤 동작 구간 구분에 사용하고, 독립 거리·각도 정확도와 직접 가시 코너 정확도는 해당 참조가 확보된 항목만 평가한다. 사람이 확인한 PnP 보조 기하 참조의 영상 오차는 별도 탐색 패널로 구분한다. 완료된 no-pose 구간의 최장은 두 방법 모두 2.535초였고 25개 저장 프레임을 포함했다. 센서 기록 공백은 이 값에 포함하지 않았으며, 저장되지 않은 프레임의 상태를 추정하지 않았다.

세션 경계를 제외하고 양쪽 자세가 모두 유효한 인접 8,737쌍에서 Base와 N3의 $|\Delta x|$ 중앙값은 각각 0.094와 0.100cm, $|\Delta z|$ 중앙값은 0.298와 0.320cm, $|\Delta\psi|$ 중앙값은 0.108과 0.114도였다. P90은 $x$와 $z$에서 각각 0.609→0.593cm와 2.150→2.059cm로 낮아졌고 방향각은 0.431→0.436도로 높아져 변화가 혼합되어 있다. 이 값은 이동과 정지를 함께 포함한 인접 출력 변화이며 정지 잡음이나 정확도가 아니다. 별도 12장의 주석은 버전별로 보존하였다. 원래 G 저장본의 기하 참조는 직접 클릭 66점과 PnP 투영 30점이며, 이후 가시성 재검수본은 직접 클릭 좌표 72점과 자체 가림 상태 24개다. 두 버전의 좌표를 섞지 않았다. 이번 별도 기하 참조 패널에는 원래 G 저장본의 96좌표만 사용했고, 사람이 고정 선택 박스를 확인하여 12장 모두 같은 대상으로 판정하였다. 공식 120장·24장 반복 검수의 직접 가시 코너 계약은 이 별도 패널로 완료 처리하지 않았다. 확인된 정지 구간과 독립 물리 참조도 없어 관련 값은 \notrun 으로 유지한다.

원래 참조와 고정 원시 예측을 대조한 표~\ref{tab:lifter_assisted_geometry}에서는 중앙 오차가 3.974→4.184px로 악화되었지만 P90은 9.643→8.977px, PCK는 88/96→91/96(91.667→94.792\%)으로 개선되어 결과가 혼합되었다. 모든 점과 네 세션을 유지했으며 오차를 보고 표본·seed를 고르지 않았다. 기하 보조 참조의 생성 시점에 예측 노출을 독립 통제했는지는 확인되지 않았고, 이번 대상 대응 단계에서 선택 박스를 보았다는 사실은 기록했다. 따라서 이 사례로 독립 블라인드 코너 복원, 물리 자세 정확도, 추적 정확도, 포크 삽입 성공률 또는 정렬시간 개선을 주장하지 않는다.
\begin{table}[t]\centering\footnotesize
\caption{사람이 확인한 PnP 보조 기하 참조의 별도 12장 탐색 패널. YOLO Base와 N3(seed 1).}\label{tab:lifter_assisted_geometry}
\begin{tabularx}{\columnwidth}{Yrrrrr}\toprule
방법 & 중앙(px) & P90(px) & PCK(\%) & 적중 & 유효/실패 \\\midrule
Base & 3.974 & 9.643 & 91.667 & 88/96 & 96/0 \\
N3 & 4.184 & 8.977 & 94.792 & 91/96 & 96/0 \\
\bottomrule\end{tabularx}
\tabnote{4세션·12장·96개 기하 코너의 고정 분모다. 원래 직접 클릭 66점과 G로 채운 PnP 투영 30점을 같은 규칙으로 사용한다. PCK는 오차 $\leq10$px이며 잘못된 대상·결측은 실패로 유지한다. 중앙값/P90은 유효 일치점 조건부 통계다. 공식 120장/24장 가시 코너 계약이나 독립 물리 T/R 정확도에 해당하지 않는다.}
\end{table}


\begin{table}[t]
\centering\footnotesize
\caption{네 기록 영상 8,910프레임의 고정 모델 오프라인 출력.}\label{tab:case}
\begin{tabularx}{\columnwidth}{Ycc}\toprule
항목 & Base & Base+N3 \\\midrule
평가 영상 / 저장 프레임 & 4 / 8,910 & 4 / 8,910 \\
가시 코너 중앙값 / P90 (px) & \notrun/\notrun & \notrun/\notrun \\
새 자세 산출률 (\%) & 98.451 & 98.451 \\
최장 완료 no-pose 구간 (초) & 2.535 & 2.535 \\
인접 $|\Delta x|$ 중앙/P90 (cm) & 0.094 / 0.609 & 0.100 / 0.593 \\
인접 $|\Delta z|$ 중앙/P90 (cm) & 0.298 / 2.150 & 0.320 / 2.059 \\
인접 $|\Delta\psi|$ 중앙/P90 (도) & 0.108 / 0.431 & 0.114 / 0.436 \\
정지 구간 전방 변동 (cm) & \notrun & \notrun \\
정지 구간 횡방향 변동 (cm) & \notrun & \notrun \\
정지 구간 방향각 변동 (도) & \notrun & \notrun \\
전방 거리 참조 오차 (cm) & \notrun & \notrun \\
횡방향 참조 오차 (cm) & \notrun & \notrun \\
방향각 참조 오차 (도) & \notrun & \notrun \\
\bottomrule\end{tabularx}
\tabnote{두 방법의 fresh/held/no-pose는 각각 8,772/0/138이다. 인접 변화는 양쪽 출력이 유효한 같은 8,737쌍의 절댓값이며 이동과 정지를 함께 포함한다. 변동은 확인된 정지 구간의 표준편차를 뜻하므로 사람 확인 전까지 x다. 이 표의 가시 코너 정확도는 공식 120장/24장 참조 계약에 해당하며 x를 유지한다. 별도 12장 PnP 보조 기하 참조 오차는 표~\ref{tab:lifter_assisted_geometry}에 구분한다. 독립 물리 정확도는 x다.}
\end{table}
